% Chapitre « Composition des équipes, horizon H2 » (partie 6, sous-étape 6.14, jalon J9).
% Sources : ADR-0013, ADR-0047 à ADR-0052 ; code de src/foot_predictor/features/ (lineup_stability.py,
% xi_quality.py, xi_quality_estimation.py, coach.py), de src/foot_predictor/modeling/ (design.py,
% models/regularized.py, models/boosting.py). Les chiffres viennent de docs/resultats/ (h2_ablations.md,
% h2_structures.md, test_scelle_h2.md), jeu ds-2026-11-01-e0defb32 pour la validation, 4 plis, 7 156 matchs.
% En cas d'écart, le code et les ADR font foi.

\chapter{Composition des équipes (horizon H2)}
\label{chap:composition}

L'horizon H1 s'arrête avant la publication des compositions. L'horizon H2 se place une heure avant le
coup d'envoi (\og T-60 \fg{}) : le onze de départ de chaque équipe et son entraîneur sont connus. Ce chapitre
décrit les trois groupes de variables que cela permet (G6 stabilité du onze, G5 qualité du onze, G4
entraîneur), le modèle H2 qui en est issu et les résultats qui l'ont déterminé. H2 est un horizon de
\emph{rejeu} : après la fin de l'abonnement à l'API, aucune composition n'arrive plus, et l'application
répond \og indisponible \fg{} en direct (ADR-0010, ADR-0047).

% ---------------------------------------------------------------------------------------------
\section{Règle temporelle avec composition}
\label{sec:composition-regle}

Pour une équipe \(e\) au match \(m\) du jour \(J\), notons \(X_m \subset \mathcal{J}\) son onze annoncé
(onze joueurs identifiés). Les variables de H2 sont des fonctions de \(X_m\), de l'entraîneur annoncé et
de l'\emph{historique du club} \(\mathcal{K}_{e,<J}\) : ses matchs de championnat des échelles, terminés,
d'un jour \emph{strictement antérieur} à \(J\), dont le onze est complet (onze titulaires identifiés, sans
entrée exclue pour collision, ADR-0020). Le match \(m\) lui-même n'y figure jamais : sinon le recouvrement
du onze avec lui-même vaudrait 1. Les calculs traversent les saisons, sans remise à zéro, et se limitent au
club. Une valeur incalculable reste \emph{vide}, jamais remplacée. L'invariance à la date de coupe
(section~\ref{sec:regle}) est vérifiée sur les 92 colonnes du jeu, au bit près.

% ---------------------------------------------------------------------------------------------
\section{Stabilité du onze (groupe G6)}
\label{sec:g6}

Soit \(\mathcal{K}_{e,<J} = (k_1, k_2, \dots)\) l'historique du club, du plus récent au plus ancien, et
\(s_k = |X_m \cap X_k|\) le nombre de titulaires que le onze de \(m\) partage avec celui du match \(k\)
(0 à 11).

\paragraph{Recouvrement.} \code{xi\_overlap\_prev} \(= s_{k_1}\), avec le match précédent du club.

\paragraph{Cohésion par paires.} Avec un poids \(w_r = 2^{-(r - 1)/h}\) pour le \(r\)-ième match en
remontant (le précédent pèse 1) et une demi-vie \(h \in \{3, 5, 10, 20\}\) matchs (\(h = \infty\) : poids
égaux),
\[
  \mathrm{coh}_h(m) = \frac{\sum_{r \ge 1} w_r\, \dfrac{s_{k_r}(s_{k_r} - 1)}{110}}{\sum_{r \ge 1} w_r}
  \in [0, 1].
\]
Le dénominateur \(110 = 11 \times 10\) est le nombre de paires ordonnées de deux joueurs distincts d'un
onze. L'expression est la somme, restreinte aux paires de \(X_m\), de la matrice \(C\) des fréquences
pondérées avec lesquelles deux joueurs ont débuté ensemble. En effet, si \(x_m\) et \(x_k\) sont les
vecteurs indicateurs des onze, le nombre de paires ordonnées \((i, j)\), \(i \ne j\), de \(X_m\) qui sont
aussi toutes deux dans \(X_k\) vaut
\[
  \sum_{i \ne j} x_{m,i}\, x_{m,j}\, x_{k,i}\, x_{k,j}
  = \Big(\sum_i x_{m,i} x_{k,i}\Big)^2 - \sum_i x_{m,i} x_{k,i}
  = s_k^2 - s_k = s_k(s_k - 1).
\]
Cette égalité, vérifiée par un test sur des matrices construites explicitement, évite de construire la
matrice : le calcul est linéaire en la taille de l'historique. Une cohésion de 1 veut dire que le onze a
toujours joué ensemble. La décroissance se fait par \emph{rang} de match, et non en jours : la pause d'été
ne fait pas vieillir l'historique.

\paragraph{Nouveaux titulaires.} \code{xi\_new\_starters} est le nombre de joueurs de \(X_m\) qui ont moins
de 3 titularisations de championnat dans le club sur les 365 jours précédents.

\paragraph{Variantes de diagnostic.} La spécification d'origine remettait la cohésion à zéro à chaque saison
(\code{xi\_cohesion\_reset}, vide au premier match) ; une autre prend les 10 derniers matchs à poids égaux
(\code{xi\_cohesion\_win10}). Elles sont mesurées, jamais choisies sur les plis de test.

% ---------------------------------------------------------------------------------------------
\section{Qualité du onze (groupe G5)}
\label{sec:g5}

\paragraph{Note rétrécie d'un joueur.} Une apparition de \(w\) minutes (\(1 \le w \le 130\)) donne une note
\(r\) (retenue entre 3 et 10, sinon l'apparition n'a pas de note). On modélise, pour un joueur \(j\) de poste
\(p\),
\[
  r_{j,t} = \theta_j + \varepsilon_{j,t},
  \qquad \operatorname{Var}(\varepsilon_{j,t}) = \frac{\sigma^2}{w_{j,t}},
  \qquad \theta_j \sim (\mu_p, \tau^2).
\]
Pour un joueur de \(M\) minutes notées avant \(J\) et de moyenne pondérée par les minutes \(\bar r_j\), la
meilleure prévision linéaire de \(\theta_j\) est la note \emph{rétrécie}
\[
  \tilde r_j = \frac{M\, \bar r_j + k\, \mu_p}{M + k},
  \qquad k = \frac{\sigma^2}{\tau^2},
\]
où \(\mu_p\) est la moyenne, pondérée par les minutes, des notes antérieures à \(J\) au poste \(p\). Comme
pour les moyennes glissantes (section~\ref{sec:glissants}), c'est une combinaison convexe : un joueur sans
note antérieure vaut \(\mu_p\), et ces cas sont comptés (ADR-0013, règle 3 ; ADR-0047, D9). La constante
\(k\), en minutes, est estimée \emph{une seule fois} par la méthode des moments (analyse de variance
déséquilibrée), sur 2015-16 à 2019-20 seulement, donc dans l'apprentissage de tous les plis :
\(\hat\sigma^2 \approx 20{,}6\), \(\hat\tau^2 \approx 0{,}027\), \(k = 772{,}6\) minutes (9\,001 joueurs,
446\,381 apparitions). Environ neuf matchs complets suffisent donc à donner autant de poids à la moyenne
du joueur qu'à celle de son poste. Le poste est celui de la composition annoncée, à défaut le dernier
connu ; le poste d'un match passé ne vient jamais des statistiques de ce match pour rétrécir.

\paragraph{Les trois variables.} Pour l'équipe :
\begin{itemize}
  \item \code{xi\_quality} \(= \frac{1}{11}\sum_{j \in X_m} \tilde r_j\) ;
  \item \code{xi\_regulars\_missing} \(= |\mathcal{R} \setminus X_m| \,/\, |\mathcal{R}|\), où
        \(\mathcal{R}\) est l'ensemble des \emph{habituels} du club : les joueurs titulaires d'au moins 5 des
        10 derniers matchs de l'historique (vide avec moins de 10 matchs ou sans habituel). Un départ compte
        comme une absence tant que le joueur reste dans la fenêtre ;
  \item \code{xi\_mean\_age} : âge moyen des titulaires en années au jour de \(m\), vide si une date de
        naissance manque.
\end{itemize}
Les habituels absents sont corrélés à \(-0{,}85\) avec la cohésion de G6, et \code{xi\_quality} à \(0{,}83\) avec
l'Elo : à Elo donné, la corrélation tombe à \(0{,}07\). La qualité du onze en grande partie \emph{redit} la
force de l'équipe, d'où un gain attendu petit.

% ---------------------------------------------------------------------------------------------
\section{Entraîneur (groupe G4)}
\label{sec:g4}

La \emph{période en cours} d'un entraîneur \(c\) est la suite ininterrompue des derniers matchs de
l'historique du club dirigés par \(c\) ; elle est vide si le dernier match avait un autre entraîneur (\(c\)
arrive) ou si l'entraîneur est inconnu. Alors \code{coach\_tenure\_days} est le nombre de jours entre le premier
match de la période et \(J\), plafonné à 365 (0 pour un entraîneur qui arrive), et \code{coach\_is\_new}
vaut 1 si la période compte moins de 5 matchs. L'API ne donne l'entraîneur qu'à partir de 2016-17 : un
entraîneur déjà en poste en août 2016 paraît arrivé à cette date (\emph{censure à gauche}, 19,5~\% de \og
nouveaux \fg{} contre 10~\% ensuite), d'où un apprentissage de G4 à partir de 2017-18 (ADR-0047, D3).

% ---------------------------------------------------------------------------------------------
\section{Le modèle H2}
\label{sec:modele-h2}

Le modèle H2 garde la structure M3 (section~\ref{sec:m3}) et ajoute au prédicteur linéaire les variables
du groupe retenu, pour l'équipe \(e\) et pour son adversaire, standardisées sur l'apprentissage du pli. Avec
\(g_e = (\code{xi\_quality}, \code{xi\_regulars\_missing}, \code{xi\_mean\_age})\) de l'équipe,
\[
  \log \lambda_{m,e} = \underbrace{\beta_0 + \beta_{\mathrm{ch}(m)} + \cdots + \sum_{j=1}^{8} \gamma_j\,
  \widetilde{\log g_j}}_{\text{modèle H1 (section~\ref{sec:retenu})}}
  + \delta^\top \tilde g_e + \delta'^\top \tilde g_{\mathrm{adv}(e)},
  \qquad Y_{m,e} \sim \mathcal{P}(\lambda_{m,e}),
\]
soit six coefficients de plus. Comme les notes de 2015-16 n'ont pas d'historique pour rétrécir
(\emph{rodage}, ADR-0047, D3), l'apprentissage de G5 commence en 2016-17, et les lignes de 2015-16 sont
comptées, non utilisées. Les lignes dont une variable manque sont exclues et comptées (0,46~\% des lignes
des plis), jamais remplies. Les comparaisons portent sur l'intersection des matchs.

% ---------------------------------------------------------------------------------------------
\section{Décisions de la validation glissante}
\label{sec:h2-decisions}

La règle de décision de l'ADR-0037 s'applique \emph{mécaniquement}, dans l'ordre G6, G5, G4 (chaque
groupe s'ajoutant au modèle retenu juste avant lui), sur la spécification écrite avant toute mesure
(ADR-0048). Gain de log-loss du total sur 7\,156 matchs (positif : le groupe aide) :

\begin{center}
\begin{tabular}{llccl}
\toprule
Étape & Comparaison & Gain [IC 95~\%] & Plis gagnés & Verdict \\
\midrule
1 & base H1 \(\to\) + G6 & \(+0{,}00035\) \([-0{,}00022 ; +0{,}00091]\) & 3/4 & non retenu \\
2 & base H1 \(\to\) + G5 & \(+0{,}00078\) \([+0{,}00003 ; +0{,}00151]\) & 4/4 & \textbf{retenu} \\
3 & G5 \(\to\) + G4 & \(-0{,}00090\) \([-0{,}00158 ; -0{,}00022]\) & 1/4 & non retenu \\
\bottomrule
\end{tabular}
\end{center}

Le modèle H2 est donc M3 sur G0, G1, G2 et G5. Le gain de G5 est minuscule (0,04~\% du log-loss) et passe
la règle de justesse : sa borne basse est à \(+0{,}00003\) (Diebold-Mariano \(p = 0{,}040\)). À fenêtre
d'apprentissage égale (dès 2016-17), il vaut \(+0{,}00061\) \([-0{,}00010 ; +0{,}00130]\) : une partie du gain
tient à la fenêtre. Il se concentre aux journées 1 à 5 (\(+0{,}0021\) \([+0{,}0002 ; +0{,}0041]\)), quand les
moyennes glissantes sont encore peu informatives. La stabilité (G6) est \emph{abandonnée} pour H2 : aucune
des huit variantes (cohésion à demi-vie 3, 5, 10, 20 matchs ou infinie, remise à zéro par saison, fenêtre de
10 matchs, recouvrement seul) ne se distingue de la base.

\paragraph{Structures (ADR-0049).} Avec les mêmes variables, le ridge (M6, chapitre~\ref{chap:modeles})
ne remplace pas M3 (\(+0{,}00014\), \([-0{,}00044 ; +0{,}00072]\), 3 plis sur 4). Le boosting à perte de
Poisson (M7), \code{HistGradientBoostingRegressor(loss="poisson")}, minimise la déviance de Poisson
\(2\sum_i [y_i \log(y_i/\lambda_i) - (y_i - \lambda_i)]\) sur un ensemble d'arbres de profondeur 2 à 4, et
perd significativement (\(-0{,}00235\), \([-0{,}00444 ; -0{,}00025]\), 0 pli sur 4) : sur un signal aussi
petit, une structure flexible ajuste du bruit. La structure reste M3.

% ---------------------------------------------------------------------------------------------
\section{Test scellé de l'horizon H2}
\label{sec:test-scelle-h2}

Les mêmes 2\,092 matchs scellés que pour H1 (du 1\textsuperscript{er} juillet 2025 au 18 octobre 2026) sont
utilisés une dernière fois (ADR-0047, ADR-0050) : apprentissage sur 2015-16 à 2024-25 (G5 dès 2016-17),
hyperparamètres par validation interne, réajustement, une seule évaluation. La liste, figée avant le test
(tag \code{pre-scelle-h2}), compare \(M_{\mathrm{H1}}\) (M3 sur G0, G1, G2) à \(M_{\mathrm{H2}}\) (M3 sur
G0, G1, G2, G5). Avant l'évaluation, les 68 colonnes H1 du jeu scellé de H2 ont été vérifiées
\emph{identiques au bit près} à celles du jeu du premier test. 169 matchs n'ont pas de valeur de G5
(surtout l'âge moyen, faute de date de naissance) et sont écartés : \(1\,923\) matchs sont évalués.

\begin{center}
\begin{tabular}{lcccc}
\toprule
Modèle & Log-loss & RPS & Brier \(P(T > 2{,}5)\) & Pente de calibration \\
\midrule
\(M_{\mathrm{H1}}\) & 1,8741 & 0,1270 & 0,2419 & 0,97 \\
\(M_{\mathrm{H2}}\) & 1,8728 & 0,1268 & 0,2415 & 0,99 \\
B1 & 1,8905 & 0,1290 & 0,2457 & 0,93 \\
\bottomrule
\end{tabular}
\end{center}

La comparaison de décision, \(M_{\mathrm{H1}} - M_{\mathrm{H2}}\) en log-loss du total, vaut
\(+0{,}00129\) \([-0{,}00035 ; +0{,}00289]\) (Diebold-Mariano \(p = 0{,}12\)) : même signe et même ordre de
grandeur que la validation glissante, \emph{non significatif}. C'est le résultat attendu : l'ADR-0047
estimait qu'un gain inférieur à environ \(0{,}007\) ne se distinguerait probablement pas de zéro sur ces
matchs, et celui de G5 en est près de neuf fois plus petit. Ce résultat ne retire donc pas G5, et ne le
confirme pas au sens statistique ; la conclusion principale reste celle de la validation glissante. Le
modèle H2 bat B1 de \(+0{,}0177\) \([+0{,}0089 ; +0{,}0268]\). Les prédictions de \(M_{\mathrm{H1}}\) sont
identiques au bit près à celles de la relance du test de H1.
